\section{Future work}
\label{sec:future}
The following architectures are proposed and have not been implemented. Each study pairs direct and scalar experts with matched initialization, data, and update exposure. Trajectory supervision is the first planned experiment.

\subsection{Trajectory supervision}
\label{sec:trajectory}
An auxiliary head reconstructs two cubic Hermite segments for each of the six arm controls over a ten-action chunk. For fixed segment intervals,
\begin{equation}
q(s)=(2s^3-3s^2+1)p_0+(s^3-2s^2+s)m_0+(-2s^3+3s^2)p_1+(s^3-s^2)m_1,
\quad s\in[0,1].
\end{equation}
The head maps pooled action-expert features to endpoints and segment-scaled tangents. Decoding these coefficients at the command times gives the training objective
\begin{equation}
\mathcal L=\mathcal L_{\rm flow}+10\,\operatorname{MSE}(\widehat A_{\rm arm},A_{\rm arm}).
\end{equation}
Gripper supervision remains unchanged, and the auxiliary head is removed at deployment. The proposed biased $1024\!\to\!512\!\to\!512\!\to\!48$ MLP adds 812,080 parameters, about 0.19\% of the trainable parameters \cite{lv2026hermitecode}.

Hermite-VLA reports an overall improvement from 95.9\% to 98.7\%, while its one-segment variant reaches 93.7\% \cite{lv2026hermite}. The planned factorial comparison crosses direct/scalar experts with no auxiliary loss, Hermite reconstruction, and an equally sized cosine-basis head. Feature access, coefficient counts, and loss scales are matched. Measurements include held-out reconstruction error, success, gripper timing, segment-join error, and replanning discontinuities. The direct/scalar difference within each supervision condition tests the scalar contribution.

\subsection{Stage prediction}
A stage predictor maps images, the instruction, current robot state, and a short causal state history to the next gripper event, its end-effector pose, and time remaining. Demonstration labels follow a fixed extraction rule; terminal poses carry an explicit no-further-event label. The predicted waypoint conditions several expert layers.

Waypoint hierarchies and $\pi_{0.7}$ motivate this design \cite{xie2026fastplans,physicalintelligence2026pi07}. The comparison includes state alone, state with history, and state with history and predicted waypoints, with identical inputs for direct and scalar models. Erased or swapped waypoints on development resets test whether execution depends on the prediction. Simulator outcomes for first-object placement, second-object acquisition, and final retention distinguish stages that existing timeout records cannot resolve.

\subsection{Geometric invariance}
Arm commands use calibrated gripper-relative coordinates, with equal normalization scales across axes mixed by tabletop yaw and an invariant gripper command. Let $R$ be an orthogonal action representation and $T_R$ its context transform. The proposed potential satisfies
\begin{equation}
\Psi_\theta(Rx,\tau,T_Rc)=\Psi_\theta(x,\tau,c)
\quad\Longrightarrow\quad
\nabla_{Rx}\Psi_\theta(Rx,\tau,T_Rc)=R\nabla_x\Psi_\theta(x,\tau,c).
\end{equation}
The implication requires differentiability, orthogonality of $R$, and invariance of every action-dependent path. An invariant network or finite-group averaging can enforce the premise; averaging gives exact symmetry only for that group. The coordinate action must include geometric context; image features require a defined transformation law. Passive coordinate changes and active scene rotations require separate tests. Distributional equivariance also requires compatible noise and sampling.

AnyCamVLA's multi-view control improves novel-view success from 69.9\% to 87.2\%, motivating a strong augmentation baseline \cite{heo2026anycam}. Following data-efficiency studies such as SAVLA \cite{li2026savla}, the comparison includes augmentation, relative-coordinate direct, equivariant direct, and invariant scalar models. Calibrated rerenders share source-trajectory splits. Outcomes comprise numerical equivariance error and success across demonstration budgets and held-out viewpoints.

\subsection{Future prediction}
Given policy features and five demonstrated actions, a training-only predictor estimates the target
\begin{equation}
\Delta e_t=\operatorname{stopgrad}\bigl(e(I_{t+5})-e(I_t)\bigr),
\end{equation}
where $e$ is a frozen public image encoder. Predictor gradients update the shared policy representation. Image/action windows remain within episodes, and all controls receive the same backbone updates.

DYNA-2 motivates future supervision, while SaiVLA-JEPA finds no detectable aggregate closed-loop gain after matching current features \cite{dyna2,saivlajepa}. Controls comprise current-frame, action-conditioned future-change, and within-task shuffled-future supervision at matched capacity and cost. Swapped action prefixes from identical simulator states test action sensitivity. The primary outcome is closed-loop success relative to the current-frame control.

\subsection{Convex generation}
\label{sec:convex}
For $x_\tau=(1-\tau)A+\tau\epsilon$, ideal endpoint regression gives
\begin{equation}
v^*(z,1,c)=z-\mathbb E[A\mid c],\qquad z-v^*(z,1,c)=\mathbb E[A\mid c].
\end{equation}
This assumes independent noise and an endpoint value or regular extension to $\tau=1$, which continuous-time training samples with zero probability. It applies to both ideal direct and scalar regression; finite learned networks need not satisfy it.

The proposed distributional model samples a categorical gripper schedule $G$ and generates arm commands through
\begin{equation}
\widehat A=\nabla_z\Phi_\theta(z,c,G),\qquad z\sim\mathcal N(0,I_{60}),\qquad
\Phi_\theta(\cdot,c,G)\text{ convex}.
\end{equation}
An input-convex network uses nonnegative hidden-to-hidden weights, convex nondecreasing activations, and affine noise skips \cite{icnn}. Conditioning must preserve convexity in noise. Schedules receive categorical log loss. Arms, conditioned on demonstrated schedules, receive
\begin{equation}
\mathcal L_{\rm ES}=\mathbb E\lVert\widehat A-A\rVert_2
-\tfrac12\mathbb E\lVert\widehat A-\widehat A'\rVert_2,
\end{equation}
using generated draws independent of each other and the target conditional on $(c,G)$. Schedule support retains observed multi-switch chunks.

The energy score is strictly proper for distributions with finite first moments. If the model family contains the target conditional distribution, that distribution uniquely minimizes the population loss \cite{gneitingraftery2007}. If a convex gradient transports Gaussian noise to a target with finite second moments, it is the almost-everywhere unique quadratic-cost optimal transport map \cite{brenier1991,mccann1995}. Finite-network approximation and optimization remain empirical questions; singular targets may require nonsmooth potentials.

ESP and DriftingVLA provide one-step baselines \cite{esp,driftingvla}. An unrestricted generator provides the matched control, sharing the gripper head, loss, noise draws, and budget. Gaussian tests assess distribution recovery for both models and the convex model's agreement with the analytic transport map. Robot tests measure success, calibration, and latency including the input derivative.

\subsection{Evaluation}
\label{sec:confirmation}
The evaluation will compare full-model direct/scalar pairs initialized from \texttt{pi05\_base}, with matched inputs, losses, and backbone updates. A separate \texttt{pi05\_libero} rerun uses its own normalization. This comparison requires conditioning-backbone gradients in the trainer. The protocol retains ten predicted actions, five executed actions, original task limits, paired resets/noise, and one qualified hardware family. Flow policies use $K=1,10$; endpoint generators use one call. We will report success, GPU-hours, memory, and inference cost.

Selection uses development resets, followed by the official $40\times50$ starts and confirmation of one finalist on sealed fresh resets. The planned budget is ten paired training seeds and 100 fresh starts per task, subject to pilot variance. Fresh starts require a pinned independent sampler: the release provides only 50 fixed starts. Sealed-state hashes and validity rules are fixed before policy comparison. Paired intervals cover training seeds and shared starts over the fixed 40-task benchmark.

Let $\widehat\Delta$ be the estimated gain over the matched direct model, in percentage points, and $L_{.95}$ the lower endpoint of its paired 95\% confidence interval. The proposed criterion is $\widehat\Delta\geq1$ with $L_{.95}>0$. A claim that the true gain exceeds one point requires $L_{.95}>1$. Corrected LIBERO-Plus provides the planned robustness test \cite{liberoplus,lajepa}.

Nested budgets contain 5, 10, 20, and all demonstrations per task: 200, 400, 800, and 1,693 total. Models share whole-episode splits and subset-fitted normalization. A 400-demo candidate must beat its matched direct control. If its lower confidence bound against the 1,693-demo direct model also exceeds the predeclared $-1$-point noninferiority margin, it supports comparable success with 4.23 times fewer downstream demonstrations. This claim concerns adaptation from the shared pretrained base.
